\documentclass[12pt]{article}
\usepackage[utf8]{inputenc}
\usepackage[russian]{babel}
\usepackage{amsmath}
\usepackage{amssymb}
\usepackage{geometry}
\geometry{a4paper, margin=2cm}

\begin{document}

\section*{Задача 1.4}

\textbf{Условие:} Сформулировать программу прибавления $k$
к двоичному числу, записанному наоборот. Провести анализ
сложности в среднем и в худшем случаях.

\section*{Программа}

Пусть $k\geq 1$ - фиксированное число, $L$ - длина
его двоичной записи, а $k[i]$ - его $i$-й бит,
где $i=0,\ldots,L-1$. Перенос обозначим через
$c\in\{0,1\}$.

\begin{verbatim}
// Number k is fixed. Let k[i] be the i-th bit of k
// The program adds k to the input number

START:
    # -> #R Add_0_0

Add_i_carry:
    0 -> {(0+k[i]+carry)%2} R Add_{i+1}_{(0+k[i]+carry)//2}
    1 -> {(1+k[i]+carry)%2} R Add_{i+1}_{(1+k[i]+carry)//2}

    // Here i=L-1
    0 -> {(0+k[i]+carry)%2} R {STOP || Inc}
    1 -> {(1+k[i]+carry)%2} R {STOP || Inc}

    // The input word has length less than L
    # -> {(k[i]+carry)%2} R Add_{i+1}_{(k[i]+carry)//2}
    # -> {(k[i]+carry)%2} R {STOP || WriteOne}

    // If i==0, the input word is empty
    # -> #H STOP

Inc:
    0 -> 1H STOP
    1 -> 0R Inc
    # -> 1H STOP

WriteOne:
    # -> 1H STOP
\end{verbatim}

\subsection*{Сложность}

Пусть $n$ - длина $w$, $L$ - длина $k$.

На пустом слове:
\[
T(\varnothing)=2
\]

Дальше рассматриваем непустое слово $w$.


Сначала один шаг из $START$, затем по одному
шагу для каждого символа $k$:
\[
T_0(n)=1+L
\]

Если после обработки последнего разряда перенос нулевой,
то STOP
\[
T(w)=L+1
\]

\paragraph{$carry=1$, $1\leq n\leq L$.}

После обработки $L$ символов - переход в $WriteOne$:
\[
T(w)=L+2
\]

\paragraph{$carry=1$, $n>L$.}

Пусть $w_{L+1}\ldots w_n = 1^r0^{n-L-r}$, $0\leq r\leq n-L$.
$Inc$ заменяет эти $r$ единиц на нули, затем записывает
единицу на первом символе $0$ или на символе $\#$:
\[
T_{Inc}(w)=r+1
\]

В итоге:
\[
T(w)=L+1+(r+1)=L+r+2
\]
При $r=n-L$ будет наибольшее время:
\[
T(w)=n+2
\]

\subsection*{Сложность в худшем}

Пусть $P_i$ - слова длины $i$, после
обработки которых остаётся перенос. Это худшие случаи.

\paragraph{$n=0$.}
\[
T_{\max}(0)=2
\]

\paragraph{$1\leq n\leq L$.}
\[
w':=w0^{L-n}
\]

Пусть $W_n := \{w: w'\in P_i\}$.

Для $w\in W_n$ как сказано выше:
\[
T(w)=L+2
\]
Для остальных:
\[
T(w)=L+1
\]

Значит, если $|W_n|>0$, худшие случаи - $w\in W_n$.
Если $|W_n|=0$, все слова обрабатываются за $L+1$ шаг.

\[
T_{\max}(n)=
\begin{cases}
L+2, & |W_n|>0\\
L+1, & |W_n|=0
\end{cases}
\]

\paragraph{$n\geq L+1$.}

Пусть $w=pu$, $|p|=L$, $|u|=n-L$.

Если $p\in P_i$, запускается $Inc$. 
Пусть $u = 1^r0*\dots*$.
Тогда
\[
T(w)=L+r+2
\]

Максимум достигается когда $r=n-L$, то есть
\[
w=p1^{n-L}, p\in P_i
\]

Ну и значит, слово $1^n$ - худший случай.
\[
T(1^n)=L+1+(n-L+1)=n+2
\]

Значит, сложность в худшем:
\[
T_{\max}(n)=
\begin{cases}
2, & n=0\\
L+2, & 1\leq n\leq L,\ |W_n|>0\\
L+1, & 1\leq n\leq L,\ |W_n|=0\\
n+2, & n\geq L+1
\end{cases}
\]

\subsection*{Сложность в среднем}

\paragraph{$n=0$.}
\[
T_{avg}(0)=2
\]

\paragraph{$n=1$.}
Cлова $0$ и $1$ - дополнительный переход нужен
только для $1$, если оно из $P_0$.:
\[
T_{avg}(1)
=\frac{(L+1)+(L+1|W_1|)}{2}
=L+1+\frac{|W_1|}{2}
\]

\paragraph{$2\leq n\leq L$.}
Слова: $w=*^{n-1}1$, всего их
$2^{n-1}$. Из них $|W_n|$ из $P_i$. Поэтому
\[
T_{avg}(n)
=\frac{(2^{n-1}-|W_n|)(L+1)+|W_n|(L+2)}{2^{n-1}}
=L+1+\frac{|W_n|}{2^{n-1}}
\]

\paragraph{$n\geq L+1$.}
Пусть $w=pu$, где $|p|=L$, $|u|=n-L$, $u=*\dots*1$. Для первых $L$ битов перенос
возникает на словах $P_i$. Каждый $P_i$ содержит $2^i$
слов. Поэтому всего $p$ с переносом:
\[
K=\sum_{i=0}^{L-1}k_i2^i=k
\]

Если переноса нет, $T(w)=L+1$. Если он есть, то рассмотрим $Inc$.
Пусть $u=1^r0z$, $z=*\dots*1$, $0\leq r\leq n-L-2$.
Таких слов $2^{n-L-r-2}$, и каждый требует $r+1$
дополнительных шагов. Ещё один вариант - $u=1^{n-L}$:
$n-L+1$ переходов. Поэтому переходов по суффиксам:
\[
S_{n-L}=\sum_{r=0}^{n-L-2}(r+1)2^{n-L-r-2}+(n-L+1)
\]
\[
\frac{S_{n-L}}{2^{n-L-1}}
=\sum_{r=0}^{n-L-2}\frac{r+1}{2^{r+1}}
  +\frac{n-L+1}{2^{n-L-1}}
=2-\frac{n-L+1}{2^{n-L-1}}+\frac{n-L+1}{2^{n-L-1}}=2
\]

Всего слов - $2^{n-1}$.
Для каждого из $K$ префиксов с переносом плюс $S_{n-L}$
дополнительных переходов. Значит,
\[
T_{avg}(n)
=L+1+\frac{KS_{n-L}}{2^{L+(n-L)-1}}
=L+1+\frac{k\,2^{n-L}}{2^{L+(n-L)-1}}
=L+1+\frac{k}{2^{L-1}}
\]

В итоге, сложность в среднем:
\[
T_{avg}(n)=
\begin{cases}
2, & n=0\\
L+1+\dfrac{|W_1|}{2}, & n=1\\
L+1+\dfrac{|W_n|}{2^{n-1}}, & 2\leq n\leq L\\
L+1+\dfrac{k}{2^{L-1}}, & n\geq L+1
\end{cases}
\]

\end{document}